% !TEX program = pdflatex
\documentclass[3p,11pt]{elsarticle}
\usepackage{amsmath,amssymb}
\usepackage{array}
\usepackage{booktabs}
\usepackage[tableposition=top]{caption}
\usepackage{enumitem}
\usepackage{tikz}
\usepackage{pgfplots}
\pgfplotsset{compat=1.18}
\usepackage{url}
\urlstyle{same}
\usepackage[hidelinks]{hyperref}
\emergencystretch=3em
% first-page footer: the page number alone
\makeatletter
\def\ps@pprintTitle{\let\@oddhead\@empty\let\@evenhead\@empty
  \def\@oddfoot{\hfil\thepage\hfil}\let\@evenfoot\@oddfoot}
\makeatother

% ---- notation ---------------------------------------------------------------
\newcommand{\um}{\,$\mu$m}
\newcommand{\kms}{km\,s$^{-1}$}
\newcommand{\flux}{erg\,cm$^{-2}$\,s$^{-1}$}
\newcommand{\mA}{m\AA}
\newcommand{\sci}[2]{#1\times10^{#2}}
\newcommand{\ion}[2]{#1\,\textsc{#2}}
\newcommand{\nuc}[2]{$^{#1}$#2}
\newcommand{\lxl}{$L_X/L_\mathrm{bol}$}

\begin{document}

\begin{frontmatter}

\title{Glints, nightside light, satellite belts, timing steps and occulter shapes: photometric technosignature searches in Kepler, TESS and Cepheid timing}

\author{Charles C. Norton}

\begin{abstract}
Artificial objects around other stars could show in their photometry as glints of
reflected light at an orbital period, as light from a planet's dark side, as a belt of
satellites crossing the star with its planet, as a timing signal written into a regular
clock, or as an occulter whose shape no planet has. We search the Kepler and TESS
archives for each. Folded at the periods of $816$ Kepler and
$794$ TESS planet candidates whose hosts show three or more brief flashes in $252239$
and $317274$ star-days of short-cadence photometry, the flashes gather at one phase no
more often than chance predicts, and the flashes in $2149$ TESS $20$\,s light curves of
stars within $30$\,pc recur at no period beyond chance. The nightsides of $63$
ultra-short-period rocky planets emit no light above $3\sigma$, with a median limit of
$102$\,ppm of the star's flux. Belts of satellites $3$ to $90$ planetary radii from
$461$ Kepler planets add no shoulders to their transits once each transit is timed and
measured from its own light curve, with a median $3\sigma$ limit of $117$\,ppm six radii
out. Of $1810$ Kepler planet candidates with $20$ or more
transit times, none shows a phase step held from one transit on that a sinusoidal
timing variation or chance does not explain, nor do the times of maximum of $2939$ LMC
Cepheids over $14.6$ years. Those of $1116$ Kepler and $255$ TESS rejected candidates
whose light curves fold to a stable profile are eclipsing binaries, stellar variability,
heartbeat binaries and smooth dips of constant depth.
\end{abstract}

\begin{keyword}
Technosignatures \sep Techniques: photometric \sep Planets and satellites: detection \sep
Stars: variables: Cepheids \sep Occultations
\end{keyword}

\end{frontmatter}

% =============================================================================
\section{Introduction}\label{sec:intro}
% =============================================================================

The photometric archives of Kepler~\cite{Borucki2010} and TESS~\cite{Ricker2015} hold
years of precise light curves of hundreds of thousands of stars, and several kinds of
artificial object would leave marks in them. Flat surfaces in orbit, such as mirrors or
solar panels, reflect sunlight specularly and appear as brief glints repeating at the
orbital period~\cite{Jaiswal2023}. Artificial lighting on a planet's nightside would
show as light before and after transit that the planet's own heat does not
explain~\cite{LoebTurner2012}. An occulter of non-planetary shape would give a transit
profile no limb-darkened disc produces~\cite{Arnold2005,Wright2016}. And a civilization
could signal by shifting the phase of a regular clock: the transits of a body it
controls, or the pulsations of a Cepheid it perturbs~\cite{Learned2012}.

Where these ideas have been tested, it has been on individual objects. This paper
applies each to the whole of an archive, with nulls drawn from the same light curves,
so that the number of signals expected by chance is known.

\subsection{Contributions}
\begin{enumerate}[leftmargin=*,itemsep=1pt]
\item Searches for glints at every Kepler and TESS planet-candidate period and at any
period in the TESS $20$\,s light curves of the nearest stars.
\item A uniform limit on the nightside light of $63$ ultra-short-period rocky planets.
\item A search of the transits of $461$ Kepler planets for the shoulders of satellite
belts.
\item Searches for held phase steps in the transit times of $1810$ Kepler candidates and
in the pulsation timing of $2939$ LMC Cepheids.
\item Shape tests of $1116$ Kepler and $255$ TESS rejected candidates.
\end{enumerate}

% =============================================================================
\section{Searches}\label{sec:searches}
% =============================================================================

\emph{Glints at planet periods.} For the $765$ hosts of Kepler planet candidates with
short-cadence data and the $6355$ hosts of TESS objects of interest (TOIs) that are
neither false positives nor false alarms~\cite{Guerrero2021}, each light curve, divided
by a running median of $30$ (Kepler) or $62$ (TESS) minutes, gives flashes: brightenings
five robust deviations high lasting one to three cadences, with quiet flux on either
side. For every candidate on a star with three or more flashes, the flashes are folded at
its period, and the busiest window of $2.5$ cadences in phase is compared with the same
statistic for $400$ sets of flash times drawn from the star's own good cadences. The
per-cadence scatter of the detrended light curves has a median of $780$\,ppm in the
Kepler short-cadence sample and $1600$\,ppm in the TESS two-minute sample, so a flash
must reach about $0.4$ and $0.8$ per cent of the star's flux. The quiet flux required on
either side, four cadences within two robust deviations, passes a flash in white noise
with probability $0.954^8=0.69$, and the package's end-to-end test of the detector on
synthetic Kepler cadences recovers $0.62$ of isolated flashes at ten deviations, so the
searches see about two thirds of the flashes above their thresholds. A flat mirror of area $A$
and reflectivity $\rho$ that reflects the star's disc toward the observer adds a glint
of $\rho A/\pi R_\star^2$ of the star's flux, whatever its distance from the star, so
those thresholds correspond to perfect reflectors of $6\times10^{15}$ and
$8\times10^{15}$\,m$^2$, $80\,000$ to $90\,000$\,km on a side, around a star of solar
radius. A mirror whose orientation follows its orbit sweeps its reflection across the
observer in $R_\star P/2\pi a$, $1.5$ to $6$ hours at $0.05$ to $1$\,au from a solar
radius, longer than the running median and the one-to-three-cadence window, so what the
search covers is a reflector that spins, whose beam crosses the observer in minutes and
recurs at the orbital phase where the geometry repeats.

\emph{Glints at any period.} The TESS $20$\,s light curves of Gaia DR3 stars within
$30$\,pc give flashes in the same way, over a $10$\,minute running median. For a star with
three or more flashes, every period implied by a pair of flashes, and its fractions down
to a fifth, is tested for flashes within $40$\,s of the predicted times, against the same
count for flash times drawn from the star's good cadences. The $20$\,s cadences scatter
by a median of $3000$\,ppm, so a flash must reach $1.5$ per cent of the flux, a
reflector of $3\times10^{15}$\,m$^2$ around an M dwarf of $0.35$ solar radii.

\emph{Nightside light.} A rocky planet's nightside faces us just before and after
transit and is hidden in secondary eclipse, and below about $1500$\,K it emits next to
nothing in the TESS band. Every TOI with a period below $1$\,d and a radius below
$2\,R_\oplus$ is folded over its two-minute sectors and fitted outside transit and
eclipse with a sinusoid and its first harmonic; the nightside flux is the fit at phase
zero less the mean flux in eclipse.

\emph{Satellite belts.} A belt of satellites $k$ planetary radii from a planet, such as
one at its synchronous radius~\cite{SocasNavarro2018}, crosses the star with the planet
and dims it for $kR_p/v$ before the planet's first contact and as long after its last,
where $v$ is the planet's speed across the star. The long-cadence light curves of the
$1243$ hosts of confirmed Kepler planets brighter than Kepler magnitude $15$, detrended by
a running median that skips each transit and its shoulders and with the system's other
transits masked, are folded on Holczer et al.'s measured transit times where they
exist~\cite{Holczer2016} and on the catalogue ephemeris otherwise. For $k$ from $3$,
beyond the Roche limit inside which natural rings lie, to $90$ or half the planet's Hill
radius, the flux from $18$ minutes beyond each contact to $kR_p/v$ beyond it is compared
with the flux just outside, for both shoulders together and each alone, against the same
statistic at $40$ random phases away from transit and eclipse. Every test above $z=4$ is
measured again with each transit's time and the transit's duration fitted from the light
curve, against a null whose phases are detrended with the same gap as the transit, and a
shoulder must exceed $z=5$ in both.

\emph{Phase steps in transit times.} Holczer et al.\ give the time of every transit of
$2599$ Kepler candidates~\cite{Holczer2016}. A perturbing planet moves transit times
sinusoidally, tides and precession move them smoothly, and the spacecraft's quarterly
rolls shift them at quarter boundaries; a step held from one transit on needs an
impulsive change of the orbit, which natural dynamics makes only in a close encounter
and which a civilization that moves a body, or that times an artificial occulter, would
make at will. The held step is the simplest symbol such a clock could carry, and the one
that a search without a model of the message can test. For each candidate
with $20$ or more good times, the observed minus calculated times are fitted with a
linear ephemeris, with a step held from one transit on (six or more transits each side)
and with the best sinusoid. A step counts when it beats the sinusoid by $10$ in $\chi^2$,
has a false-alarm probability below $0.01$ against $300$ permutations of the residuals,
and lies more than two days from a quarter boundary.

\emph{Phase steps in Cepheid timing.} Learned et al.\ proposed that a civilization could
advance a Cepheid's pulsation and write information into its times of
maximum~\cite{Learned2012}. Poleski's times of maximum of LMC Cepheids from MACHO
(1992--1999) and OGLE (1997--2008)~\cite{Poleski2008} give each of $2939$ stars a median
of $82$ maxima over $14.6$ years; each star's timing is fitted with a quadratic ephemeris
with offsets between sources, then with a step held from one cycle on or with a
companion's light-travel orbit of $1$ to $30$ years, and random walks of the star's own
excess scatter give the null.

\emph{Occulter shapes.} The Kepler DR25 candidates flagged only as not
transit-like~\cite{Thompson2018}, $1116$ of them, and the $255$ TESS false positives and
false alarms at $T\le11.5$ whose dispositions name no binary, blend or offset, are folded
at their periods with stellar variability removed, fitted with a symmetric trapezoid, and
measured for their departure from it, their ingress-egress asymmetry, their number of
dips and the agreement between halves of the data and between odd and even events. A
periodic, stable profile that no trapezoid describes and that has no secondary eclipse
is listed.

% =============================================================================
\section{Results}\label{sec:results}
% =============================================================================

\emph{Glints.} The Kepler light curves hold $3634$ flashes in $252239$ star-days; of the
$816$ candidates on stars with three or more, $9$ gather their flashes at one phase at a
false-alarm probability below $0.01$, against at most $8.2$ by chance, and $2$ reach the
$1/401$ floor of the draws, against $2.0$. The TESS light curves hold $8710$ flashes in
$317274$ star-days, and none of the $794$ TOIs on stars with three or more falls below a
false-alarm probability of $0.01$, against at most $7.9$ by chance. The $2149$ TESS
$20$\,s light curves, $1.13$ million star-hours, hold $19695$ flashes, $0.42$ a star-day,
mostly cosmic rays and instrumental events; in the $1232$ stars with three or more, the
best periodic family of each reaches $p\le0.01$ in $6$ stars against $12.3$ expected for
uniform $p$, and $p\le0.05$ in $38$ against $61.6$.

\emph{Nightside light.} For the $63$ planets the nightside flux divided by its error has
a median of $-0.33$ and a spread of $0.87$, and none exceeds $+3$; the highest, TOI~4567.01
at $2.9$, has an equally high level in eclipse, an offset of the eclipse level rather
than light from the planet. The median $3\sigma$ limit is $102$\,ppm of the star's flux.
That is far more than the starlight the planets receive: a planet of radius $R_p$ at
distance $a$ intercepts $(R_p/2a)^2$ of its star's light, $0.5$ to $33$\,ppm for these
$63$ (median $6.8$\,ppm), so the limits lie $3$ to $100$ times (median $21$) above all
the starlight a nightside could re-emit, and constrain only light powered by something
other than the star: at the median limit, $4\times10^{21}$\,W, with a range of
$2\times10^{20}$ to $10^{23}$\,W, which is $10^{4}$ to $10^{6}$ times humanity's
output. A nightside lit by its star's own energy, at or below a few parts per million,
needs a precision these light curves do not reach.

\emph{Satellite belts.} The $1243$ stars give $1522$ tests of $461$ planets. Eight exceed
$z=5$ with both shoulders above $z=2$, where no null phase reaches $z=5$, and every test
above $z=4$ falls below $z=5$ when measured again: with each transit's own time and fitted
duration, against the null detrended as the transit is, the largest of the $22$ planets
on the $16$ stars concerned reaches $z=3.3$ (KOI-1353.02). The excesses come from
transit-timing variations that a linear ephemeris smears across the contacts (KOI-134.01,
whose transits wander by about ten hours, falls from $z=28.7$ against the matched null to
$1.0$ on its own times, and KOI-1474.01 from $16.8$ to $-0.2$), from catalogue durations
shorter than the transit (KOI-620.01 lasts $7.4$\,h against the catalogue's $5.8$\,h), and
from the running median's gap at the transit, which the matched null reproduces
(KOI-1375.01 falls from $6.8$ to $3.0$). A shoulder six planetary radii out is excluded at
$3\sigma$ down to $117$\,ppm of the star's flux for the median planet and $22$\,ppm for the
best.

\emph{Phase steps.} Of the $1810$ Kepler candidates, $8$ steps beat the sinusoid by $10$
in $\chi^2$, one at a quarter boundary, where steps are set aside. The strongest of the
other seven follow the long-period timing variations that Holczer et al.\ list (KOI-12.01,
KOI-824.01 and KOI-2310.01), and the rest are low-signal timings at false-alarm
probabilities of $0.0012$ to $0.0020$, of which two to four are expected among the
trials. Among the Cepheids, four steps pass the tests: three fall at the start of
OGLE-II's first season (HJD $2450467$), where the data change source, and one on a date
that $15$ other stars share, so steps at a season's start or on a date other stars share
are excluded and none remains.

\emph{Occulter shapes.} Of the Kepler rejects, $129$ fold to periodic, stable profiles
at S/N $15$ or more without a secondary eclipse: $81$ are Kepler eclipsing binaries,
among them the one strongly asymmetric profile (KOI-3716.01), and the $48$ others are
rotational and pulsational variability, heartbeat binaries, or smooth V- and U-shaped
dips of constant depth (KOI-3198, KOI-8060, KOI-8240 and KOI-7698). Of the TESS rejects,
$63$ are periodic and stable over three or more events; the strongest departures are
stellar variability (TOI~1447.01), an eclipse on a blue supergiant (TOI~796.01), an
eclipse whose depth changes between sectors (TOI~7247.01) and a grazing eclipsing binary
(TOI~1090.01).

% =============================================================================
\section{Related work}\label{sec:related}
% =============================================================================

\emph{Transit anomalies.} Arnold showed how the transits of artificial occulters differ
from planets'~\cite{Arnold2005}, the $\hat{G}$ search set out the signatures of
transiting megastructures~\cite{Wright2016}, and Breakthrough Listen searched $218$
Kepler systems transit by transit for anomalies~\cite{Zuckerman2023b}.

\emph{Glints, nightsides and clocks.} Specular reflections from artificial surfaces were
proposed as a technosignature by Jaiswal~\cite{Jaiswal2023}, artificial illumination
detectable by its light curve by Loeb and Turner~\cite{LoebTurner2012}, and phase coding
of Cepheids by Learned et al.~\cite{Learned2012}, whose idea Hippke et al.\ tested on one
Kepler RR~Lyrae star~\cite{Hippke2015}. Socas-Navarro showed that a belt of satellites at
a planet's synchronous radius would add shoulders to its transits~\cite{SocasNavarro2018}.
We find no earlier archive-wide search for glints at planet periods, for nightside light
across the ultra-short-period planets, for the shoulders of satellite belts, or for held
phase steps in transit or Cepheid timing.

% =============================================================================
\section{Conclusion}\label{sec:conclusion}
% =============================================================================

No glint repeats at a planet-candidate period in Kepler or TESS, or at any period in the
fast light curves of the nearest stars, beyond what chance gives; no ultra-short-period
rocky planet shows light from its nightside above a hundredth of a per cent of its star,
twenty times what its star could power; no Kepler planet's transits carry the
shoulders of a satellite belt once each transit is timed and measured from the light
curve; no transit or Cepheid timing holds a phase
step that natural timing variations, instrument changes or chance do not explain; and
the rejected candidates with stable, non-planetary profiles are eclipsing binaries,
variable stars and smooth dips.

The glint searches apply to flashes lasting one to three cadences at five robust
deviations, $0.4$ to $1.5$ per cent of the star's flux, which perfect reflectors of
$3$ to $8\times10^{15}$\,m$^2$ would give, and to reflectors whose beam crosses the
observer in minutes; the nightside limit, $102$\,ppm, is $21$ times the starlight the
median planet intercepts and bounds only light of a non-stellar source above about
$4\times10^{21}$\,W; the belt
search applies to belts $3$ to $90$ planetary radii out whose shoulders outlast $18$ minutes
beyond contact, and the timing searches to single held steps. A glint shorter than a
cadence or longer than three, a reflector smaller than a planet, a signal at a period
other than a candidate's in the longer-cadence data, a belt
inside the Roche limit, or timing coded in many small shifts is not constrained.

\section*{Acknowledgements}

This work uses data from the Kepler and TESS missions, obtained from the Mikulski Archive
for Space Telescopes at the Space Telescope Science Institute, from ExoFOP, from the
OGLE and MACHO projects through Poleski's compilation, and from the European Space Agency
mission Gaia.

\begin{thebibliography}{99}\small\raggedright
\setlength{\itemsep}{1pt}

\bibitem{Borucki2010}
W.~J. Borucki, D.~Koch, G.~Basri, et al., Kepler planet-detection mission: introduction
and first results, \emph{Science} \textbf{327}, 977 (2010).

\bibitem{Ricker2015}
G.~R. Ricker, J.~N. Winn, R.~Vanderspek, et al., Transiting Exoplanet Survey Satellite,
\emph{J. Astron. Telesc. Instrum. Syst.} \textbf{1}, 014003 (2015).

\bibitem{Jaiswal2023}
B.~Jaiswal, Specular reflections from artificial surfaces as technosignature,
\emph{Astrobiology} \textbf{23}, 291 (2023).

\bibitem{LoebTurner2012}
A.~Loeb and E.~L. Turner, Detection technique for artificially illuminated objects in
the outer solar system and beyond, \emph{Astrobiology} \textbf{12}, 290 (2012).

\bibitem{Arnold2005}
L.~F.~A. Arnold, Transit light-curve signatures of artificial objects, \emph{Astrophys.
J.} \textbf{627}, 534 (2005).

\bibitem{Wright2016}
J.~T. Wright, K.~M.~S. Cartier, M.~Zhao, D.~Jontof-Hutter and E.~B. Ford, The $\hat{G}$
search for extraterrestrial civilizations with large energy supplies. IV. The
signatures and information content of transiting megastructures, \emph{Astrophys. J.}
\textbf{816}, 17 (2016).

\bibitem{SocasNavarro2018}
H.~Socas-Navarro, Possible photometric signatures of moderately advanced civilizations:
the Clarke exobelt, \emph{Astrophys. J.} \textbf{855}, 110 (2018).

\bibitem{Learned2012}
J.~G. Learned, R.-P. Kudritzki, S.~Pakvasa and A.~Zee, The Cepheid galactic internet,
\emph{Contemp. Phys.} \textbf{53}, 113 (2012).

\bibitem{Guerrero2021}
N.~M. Guerrero, S.~Seager, C.~X. Huang, et al., The TESS objects of interest catalog
from the TESS prime mission, \emph{Astrophys. J. Suppl.} \textbf{254}, 39 (2021).

\bibitem{Holczer2016}
T.~Holczer, T.~Mazeh, G.~Nachmani, et al., Transit timing observations from Kepler. IX.
Catalog of the full long-cadence data set, \emph{Astrophys. J. Suppl.} \textbf{225}, 9
(2016).

\bibitem{Poleski2008}
R.~Poleski, Period changes of LMC Cepheids in the OGLE and MACHO data, \emph{Acta
Astron.} \textbf{58}, 313 (2008).

\bibitem{Thompson2018}
S.~E. Thompson, J.~L. Coughlin, K.~Hoffman, et al., Planetary candidates observed by
Kepler. VIII. A fully automated catalog with measured completeness and reliability
based on data release 25, \emph{Astrophys. J. Suppl.} \textbf{235}, 38 (2018).

\bibitem{Zuckerman2023b}
A.~Zuckerman, J.~Davenport, S.~Croft, A.~Siemion and I.~de~Pater, The Breakthrough
Listen search for intelligent life: detection and characterization of anomalous
transits in Kepler lightcurves, arXiv:2312.07903 (2023).

\bibitem{Hippke2015}
M.~Hippke, J.~G. Learned, A.~Zee, et al., Pulsation period variations in the RRc Lyrae
star KIC 5520878, \emph{Astrophys. J.} \textbf{798}, 42 (2015).

\end{thebibliography}

\end{document}
